% Ngân hàng bài dummy (dữ liệu test): Hệ hai phương trình bậc nhất hai ẩn (lớp 9).
\documentclass[12pt,a4paper]{article}
\usepackage[utf8]{vietnam}
\usepackage{amsmath,amssymb}
\newenvironment{problem}[1][]{\par\medskip\noindent\textbf{Bài #1.}\ }{\par}
\newenvironment{solution}[1][]{\par\noindent\textit{Lời giải #1.}\par}{\par}

\begin{document}

% ---------------------------------------------------------------- kiểm tra cặp số có là nghiệm
\begin{problem}[9DS-HPT-032]
% nguon: dummy (tự soạn để mở rộng ngân hàng)
% lop: 9
% kc_du_kien: 9DS-HPT-01
% do_kho_du_kien: 1
Cho hệ phương trình
\[ \begin{cases} 2x - y = 4 \\ x + 3y = -5 \end{cases} \]
a) Cặp số $(1; -2)$ có là nghiệm của hệ không?

b) Cặp số $(3; 2)$ có là nghiệm của hệ không?

\begin{solution}
Một cặp số $(x_0; y_0)$ là nghiệm của hệ khi và chỉ khi nó thoả mãn đồng thời cả hai phương trình của hệ, nên ta thay từng cặp số vào cả hai phương trình.

a) Với $(1; -2)$: phương trình thứ nhất có $2 \cdot 1 - (-2) = 4$ (thoả mãn) và phương trình thứ hai có $1 + 3 \cdot (-2) = -5$ (thoả mãn). Vậy $(1; -2)$ là nghiệm của hệ.

b) Với $(3; 2)$: phương trình thứ nhất có $2 \cdot 3 - 2 = 4$ (thoả mãn), nhưng phương trình thứ hai có $3 + 3 \cdot 2 = 9 \neq -5$ (không thoả mãn). Vậy $(3; 2)$ không là nghiệm của hệ.

Đáp số: a) $\boxed{(1; -2) \text{ là nghiệm}}$; b) $\boxed{(3; 2) \text{ không là nghiệm}}$
\end{solution}
\end{problem}

% ---------------------------------------------------------------- giải hệ bằng phương pháp thế
\begin{problem}[9DS-HPT-033]
% nguon: dummy (tự soạn để mở rộng ngân hàng)
% lop: 9
% kc_du_kien: 9DS-HPT-02
% do_kho_du_kien: 1
Giải hệ phương trình
\[ \begin{cases} y = 2x - 1 \\ x + y = 8 \end{cases} \]

\begin{solution}
Phương trình thứ nhất đã cho $y$ theo $x$, nên thế $y = 2x - 1$ vào phương trình thứ hai: $x + (2x - 1) = 8$.

Thu gọn vế trái: $3x - 1 = 8$, nên $3x = 9$, suy ra $x = 3$.

Thay $x = 3$ vào $y = 2x - 1$: $y = 2 \cdot 3 - 1 = 5$.

Thử lại phương trình thứ hai: $3 + 5 = 8$ (đúng).

Đáp số: $\boxed{(x; y) = (3; 5)}$
\end{solution}
\end{problem}

% ---------------------------------------------------------------- cộng đại số, nhân một phương trình
\begin{problem}[9DS-HPT-034]
% nguon: dummy (tự soạn để mở rộng ngân hàng)
% lop: 9
% kc_du_kien: 9DS-HPT-03
% do_kho_du_kien: 2
Giải hệ phương trình
\[ \begin{cases} 2x + 3y = 7 \\ 4x - y = 7 \end{cases} \]

\begin{solution}
Nhân hai vế phương trình thứ hai với $3$ để hệ số của $y$ đối nhau với hệ số của $y$ ở phương trình thứ nhất: $12x - 3y = 21$.

Cộng vế theo vế với phương trình thứ nhất để khử $y$: $14x = 28$, suy ra $x = 2$.

Thế $x = 2$ vào phương trình thứ hai: $8 - y = 7$, suy ra $y = 1$.

Thử lại phương trình thứ nhất: $2 \cdot 2 + 3 \cdot 1 = 7$ (đúng).

Đáp số: $\boxed{(x; y) = (2; 1)}$
\end{solution}
\end{problem}

% ---------------------------------------------------------------- hệ số thập phân
\begin{problem}[9DS-HPT-035]
% nguon: dummy (tự soạn để mở rộng ngân hàng)
% lop: 9
% kc_du_kien: 9DS-HPT-03
% do_kho_du_kien: 2
Giải hệ phương trình
\[ \begin{cases} 0{,}5x + 0{,}4y = 2{,}2 \\ 0{,}3x + 0{,}2y = 1{,}2 \end{cases} \]

\begin{solution}
Nhân hai vế của mỗi phương trình với $10$ để các hệ số là số nguyên: $5x + 4y = 22$ và $3x + 2y = 12$.

Nhân hai vế phương trình thứ hai với $2$ để hệ số của $y$ bằng nhau: $6x + 4y = 24$.

Trừ vế theo vế phương trình $6x + 4y = 24$ cho phương trình $5x + 4y = 22$: $x = 2$.

Thế $x = 2$ vào $3x + 2y = 12$: $6 + 2y = 12$, nên $2y = 6$, suy ra $y = 3$.

Thử lại phương trình thứ nhất ban đầu: $0{,}5 \cdot 2 + 0{,}4 \cdot 3 = 1 + 1{,}2 = 2{,}2$ (đúng).

Đáp số: $\boxed{(x; y) = (2; 3)}$
\end{solution}
\end{problem}

% ---------------------------------------------------------------- đường thẳng đi qua một điểm và cắt trục hoành
\begin{problem}[9DS-HPT-036]
% nguon: dummy (tự soạn để mở rộng ngân hàng)
% lop: 9
% kc_du_kien: 9DS-HPT-08, 9DS-HPT-03, 8DS-HSBN-03
% do_kho_du_kien: 2
Tìm $a$ và $b$ để đường thẳng $y = ax + b$ đi qua điểm $A(2; 1)$ và cắt trục hoành tại điểm có hoành độ $5$.

\begin{solution}
Điểm cắt trục hoành có tung độ bằng $0$, nên đường thẳng đi qua điểm $B(5; 0)$.

Đường thẳng đi qua $A(2; 1)$ nên toạ độ của $A$ thoả phương trình: $1 = 2a + b$, tức là $2a + b = 1$.

Đường thẳng đi qua $B(5; 0)$ nên $0 = 5a + b$, tức là $5a + b = 0$.

Trừ vế theo vế phương trình thứ hai cho phương trình thứ nhất: $3a = -1$, suy ra $a = -\dfrac{1}{3}$; khi đó $b = -5a = \dfrac{5}{3}$.

Thử lại: với $y = -\dfrac{1}{3}x + \dfrac{5}{3}$, tại $x = 2$ có $y = -\dfrac{2}{3} + \dfrac{5}{3} = 1$ và tại $x = 5$ có $y = 0$ (đúng).

Đáp số: $\boxed{a = -\dfrac{1}{3},\ b = \dfrac{5}{3}}$
\end{solution}
\end{problem}

% ---------------------------------------------------------------- tìm hai số, phép chia có dư
\begin{problem}[9DS-HPT-037]
% nguon: dummy (tự soạn để mở rộng ngân hàng)
% lop: 9
% kc_du_kien: 9DS-HPT-07, 9DS-HPT-02
% do_kho_du_kien: 2
Tổng của hai số tự nhiên bằng $45$. Nếu lấy số lớn chia cho số nhỏ thì được thương bằng $2$ và số dư bằng $3$. Tìm hai số đó.

\begin{solution}
Gọi số lớn là $x$, số nhỏ là $y$; điều kiện $x$, $y$ là số tự nhiên và $x > y > 3$ (số dư phải nhỏ hơn số chia).

Tổng hai số bằng $45$ nên $x + y = 45$.

Số lớn chia cho số nhỏ được thương $2$ dư $3$ nên $x = 2y + 3$.

Thế $x = 2y + 3$ vào phương trình thứ nhất: $(2y + 3) + y = 45$, tức là $3y = 42$, suy ra $y = 14$; khi đó $x = 2 \cdot 14 + 3 = 31$.

Đối chiếu điều kiện: $31 > 14 > 3$ (thoả mãn). Thử lại: $31 + 14 = 45$ và $31 = 2 \cdot 14 + 3$ (đúng).

Đáp số: $\boxed{\text{Số lớn: } 31;\ \text{số nhỏ: } 14}$
\end{solution}
\end{problem}

% ---------------------------------------------------------------- tìm tham số để cặp số là nghiệm
\begin{problem}[9DS-HPT-038]
% nguon: dummy (tự soạn để mở rộng ngân hàng)
% lop: 9
% kc_du_kien: 9DS-HPT-01, 9DS-HPT-02
% do_kho_du_kien: 2
Tìm $m$ và $n$ để cặp số $(2; -1)$ là nghiệm của hệ phương trình
\[ \begin{cases} mx + ny = 4 \\ nx - my = 7 \end{cases} \]

\begin{solution}
Cặp số $(2; -1)$ là nghiệm của hệ khi và chỉ khi nó thoả mãn cả hai phương trình, nên ta thay $x = 2$, $y = -1$ vào từng phương trình.

Phương trình thứ nhất cho $2m - n = 4$; phương trình thứ hai cho $2n + m = 7$. Đây là hệ phương trình bậc nhất hai ẩn $m$, $n$.

Từ $2m - n = 4$ rút $n$ theo $m$: $n = 2m - 4$.

Thế vào $m + 2n = 7$: $m + 2(2m - 4) = 7$, tức là $5m - 8 = 7$, suy ra $m = 3$; khi đó $n = 2 \cdot 3 - 4 = 2$.

Thử lại: với $m = 3$, $n = 2$, hệ là $3x + 2y = 4$ và $2x - 3y = 7$; tại $(2; -1)$ có $6 - 2 = 4$ và $4 + 3 = 7$ (đúng).

Đáp số: $\boxed{m = 3,\ n = 2}$
\end{solution}
\end{problem}

% ---------------------------------------------------------------- đặt ẩn phụ, nghịch đảo của ẩn
\begin{problem}[9DS-HPT-039]
% nguon: dummy (tự soạn để mở rộng ngân hàng)
% lop: 9
% kc_du_kien: 9DS-HPT-06, 9DS-HPT-03
% do_kho_du_kien: 3
Giải hệ phương trình
\[ \begin{cases} \dfrac{3}{x} + \dfrac{4}{y} = -1 \\ \dfrac{6}{x} - \dfrac{2}{y} = 3 \end{cases} \]

\begin{solution}
Điều kiện xác định: $x \neq 0$, $y \neq 0$ (mẫu phải khác $0$).

Đặt $u = \dfrac{1}{x}$, $v = \dfrac{1}{y}$, hệ trở thành hệ bậc nhất hai ẩn: $3u + 4v = -1$ và $6u - 2v = 3$.

Nhân hai vế phương trình thứ hai với $2$ để hệ số của $v$ đối nhau: $12u - 4v = 6$. Cộng vế theo vế với phương trình thứ nhất: $15u = 5$, suy ra $u = \dfrac{1}{3}$.

Thế $u = \dfrac{1}{3}$ vào $3u + 4v = -1$: $1 + 4v = -1$, nên $4v = -2$, suy ra $v = -\dfrac{1}{2}$.

Trở lại ẩn ban đầu: $\dfrac{1}{x} = \dfrac{1}{3}$ nên $x = 3$; $\dfrac{1}{y} = -\dfrac{1}{2}$ nên $y = -2$.

Đối chiếu điều kiện: $x = 3 \neq 0$, $y = -2 \neq 0$ (thoả mãn). Thử lại: $\dfrac{3}{3} + \dfrac{4}{-2} = 1 - 2 = -1$ và $\dfrac{6}{3} - \dfrac{2}{-2} = 2 + 1 = 3$ (đúng).

Đáp số: $\boxed{(x; y) = (3; -2)}$
\end{solution}
\end{problem}

% ---------------------------------------------------------------- bài toán về tuổi
\begin{problem}[9DS-HPT-040]
% nguon: dummy (tự soạn để mở rộng ngân hàng)
% lop: 9
% kc_du_kien: 9DS-HPT-07, 9DS-HPT-03
% do_kho_du_kien: 3
Cách đây $5$ năm, tuổi bố gấp $7$ lần tuổi con. Sau $5$ năm nữa, tuổi bố chỉ còn gấp $3$ lần tuổi con. Tính tuổi hiện nay của mỗi người.

\begin{solution}
Gọi tuổi bố hiện nay là $x$ (tuổi), tuổi con hiện nay là $y$ (tuổi); điều kiện $x$, $y$ nguyên và $x > y > 5$ (cách đây $5$ năm người con đã có tuổi).

Cách đây $5$ năm, tuổi bố là $x - 5$ và tuổi con là $y - 5$. Bố gấp con $7$ lần nên $x - 5 = 7(y - 5)$, tức là $x - 7y = -30$.

Sau $5$ năm nữa, tuổi bố là $x + 5$ và tuổi con là $y + 5$. Bố gấp con $3$ lần nên $x + 5 = 3(y + 5)$, tức là $x - 3y = 10$.

Trừ vế theo vế phương trình $x - 3y = 10$ cho phương trình $x - 7y = -30$ để khử $x$: $4y = 40$, suy ra $y = 10$.

Thế $y = 10$ vào $x - 3y = 10$: $x = 40$.

Đối chiếu điều kiện: $40 > 10 > 5$ (thoả mãn). Thử lại: cách đây $5$ năm bố $35$ tuổi, con $5$ tuổi và $35 = 7 \cdot 5$; sau $5$ năm nữa bố $45$ tuổi, con $15$ tuổi và $45 = 3 \cdot 15$ (đúng).

Đáp số: $\boxed{\text{Bố: } 40 \text{ tuổi; con: } 10 \text{ tuổi}}$
\end{solution}
\end{problem}

% ---------------------------------------------------------------- chuyển động ngược chiều và cùng chiều
\begin{problem}[9DS-HPT-041]
% nguon: dummy (tự soạn để mở rộng ngân hàng)
% lop: 9
% kc_du_kien: 9DS-HPT-07, 9DS-HPT-03
% do_kho_du_kien: 3
Hai địa điểm $A$ và $B$ cách nhau $60$ km. Hai xe máy cùng khởi hành, một xe từ $A$ và một xe từ $B$. Nếu hai xe đi ngược chiều, hướng về nhau, thì sau $1$ giờ gặp nhau. Nếu hai xe đi cùng chiều từ $A$ về phía $B$ thì sau $3$ giờ xe xuất phát từ $A$ đuổi kịp xe xuất phát từ $B$. Tính vận tốc của mỗi xe.

\begin{solution}
Gọi vận tốc xe xuất phát từ $A$ là $x$ (km/h), vận tốc xe xuất phát từ $B$ là $y$ (km/h); điều kiện $x > y > 0$ (xe từ $A$ đuổi kịp xe từ $B$ nên đi nhanh hơn).

Đi ngược chiều, sau $1$ giờ hai xe gặp nhau, khi đó tổng quãng đường hai xe đi được bằng $AB = 60$ km: $1 \cdot (x + y) = 60$, tức là $x + y = 60$.

Đi cùng chiều, sau $3$ giờ xe từ $A$ đuổi kịp xe từ $B$, tức là xe từ $A$ đi nhiều hơn xe từ $B$ đúng $60$ km: $3x - 3y = 60$, tức là $x - y = 20$.

Cộng vế theo vế hai phương trình $x + y = 60$ và $x - y = 20$: $2x = 80$, suy ra $x = 40$; khi đó $y = 60 - 40 = 20$.

Đối chiếu điều kiện: $40 > 20 > 0$ (thoả mãn). Thử lại: $40 + 20 = 60$ và $3 \cdot 40 - 3 \cdot 20 = 60$ (đúng).

Đáp số: $\boxed{\text{Xe từ } A \text{: } 40 \text{ km/h; xe từ } B \text{: } 20 \text{ km/h}}$
\end{solution}
\end{problem}

% ---------------------------------------------------------------- bài toán lãi suất, phần trăm
\begin{problem}[9DS-HPT-042]
% nguon: dummy (tự soạn để mở rộng ngân hàng)
% lop: 9
% kc_du_kien: 9DS-HPT-07, 9DS-HPT-02
% do_kho_du_kien: 3
Một người gửi tổng cộng $800$ triệu đồng vào hai ngân hàng. Ngân hàng I trả lãi $6\%$ một năm, ngân hàng II trả lãi $8\%$ một năm. Sau một năm, tổng tiền lãi người đó nhận được là $58$ triệu đồng. Hỏi người đó đã gửi bao nhiêu tiền ở mỗi ngân hàng?

\begin{solution}
Gọi số tiền gửi ở ngân hàng I là $x$ (triệu đồng), ở ngân hàng II là $y$ (triệu đồng); điều kiện $x > 0$, $y > 0$.

Tổng số tiền gửi là $800$ triệu đồng nên $x + y = 800$.

Tiền lãi sau một năm là $6\%$ của $x$ cộng $8\%$ của $y$, bằng $58$ triệu đồng: $0{,}06x + 0{,}08y = 58$. Nhân hai vế với $100$ rồi chia cho $2$: $3x + 4y = 2\,900$.

Từ $x + y = 800$ rút $x = 800 - y$ rồi thế vào phương trình thứ hai: $3(800 - y) + 4y = 2\,900$, tức là $2\,400 + y = 2\,900$, suy ra $y = 500$.

Thay $y = 500$ vào $x = 800 - y$: $x = 300$.

Đối chiếu điều kiện: $x = 300 > 0$, $y = 500 > 0$ (thoả mãn). Thử lại: tiền lãi là $0{,}06 \cdot 300 + 0{,}08 \cdot 500 = 18 + 40 = 58$ (triệu đồng) (đúng).

Đáp số: $\boxed{\text{Ngân hàng I: } 300 \text{ triệu đồng; ngân hàng II: } 500 \text{ triệu đồng}}$
\end{solution}
\end{problem}

% ---------------------------------------------------------------- tham số, nghiệm thoả một hệ thức
\begin{problem}[9DS-HPT-043]
% nguon: dummy (tự soạn để mở rộng ngân hàng)
% lop: 9
% kc_du_kien: 9DS-HPT-05, 9DS-HPT-03
% do_kho_du_kien: 3
Cho hệ phương trình (với $m$ là tham số)
\[ \begin{cases} x + y = 2m \\ x - 2y = 3 - m \end{cases} \]
Tìm $m$ để nghiệm $(x; y)$ của hệ thoả mãn $2x - 3y = 7$.

\begin{solution}
Trừ vế theo vế phương trình thứ nhất cho phương trình thứ hai để khử $x$: $3y = 2m - (3 - m) = 3m - 3$, suy ra $y = m - 1$.

Thế $y = m - 1$ vào phương trình thứ nhất: $x = 2m - (m - 1) = m + 1$. Vậy với mọi $m$, hệ có nghiệm duy nhất $(x; y) = (m + 1; m - 1)$.

Thay nghiệm vào hệ thức $2x - 3y = 7$: $2(m + 1) - 3(m - 1) = 7$.

Bỏ ngoặc và thu gọn: $2m + 2 - 3m + 3 = 7$, tức là $5 - m = 7$, suy ra $m = -2$.

Thử lại: với $m = -2$, nghiệm là $(-1; -3)$; có $x + y = -4 = 2m$, $x - 2y = -1 + 6 = 5 = 3 - m$ và $2x - 3y = -2 + 9 = 7$ (đúng).

Đáp số: $\boxed{m = -2}$
\end{solution}
\end{problem}

% ---------------------------------------------------------------- ba đường thẳng đồng quy
\begin{problem}[9DS-HPT-044]
% nguon: dummy (tự soạn để mở rộng ngân hàng)
% lop: 9
% kc_du_kien: 9DS-HPT-01, 9DS-HPT-03, 8DS-HSBN-06
% do_kho_du_kien: 3
Cho ba đường thẳng $(d_1)\colon 2x + y = 5$, $(d_2)\colon x - y = 1$ và $(d_3)\colon mx + 2y = 6$ (với $m$ là tham số). Tìm $m$ để ba đường thẳng đồng quy, tức là cùng đi qua một điểm.

\begin{solution}
Toạ độ giao điểm của $(d_1)$ và $(d_2)$ là nghiệm của hệ $2x + y = 5$ và $x - y = 1$.

Cộng vế theo vế hai phương trình để khử $y$: $3x = 6$, suy ra $x = 2$; từ $x - y = 1$ có $y = x - 1 = 1$. Vậy $(d_1)$ và $(d_2)$ cắt nhau tại $M(2; 1)$.

Ba đường thẳng đồng quy khi và chỉ khi $(d_3)$ đi qua $M(2; 1)$, tức là toạ độ của $M$ thoả phương trình của $(d_3)$.

Thay $x = 2$, $y = 1$ vào $mx + 2y = 6$: $2m + 2 = 6$, suy ra $m = 2$.

Thử lại: với $m = 2$, đường thẳng $(d_3)$ là $2x + 2y = 6$, và tại $M(2; 1)$ có $4 + 2 = 6$ (đúng).

Đáp số: $\boxed{m = 2}$
\end{solution}
\end{problem}

% ---------------------------------------------------------------- tham số, nghiệm nằm ở một góc phần tư
\begin{problem}[9DS-HPT-045]
% nguon: dummy (tự soạn để mở rộng ngân hàng)
% lop: 9
% kc_du_kien: 9DS-HPT-05, 9DS-HPT-02, 9DS-BPT-02, 9DS-BPT-08
% do_kho_du_kien: 4
Cho hệ phương trình (với $m$ là tham số)
\[ \begin{cases} x + 2y = m \\ 2x - y = 3 \end{cases} \]
Tìm $m$ để hệ có nghiệm $(x; y)$ thoả mãn $x > 0$ và $y < 0$.

\begin{solution}
Từ phương trình thứ hai rút $y$ theo $x$: $y = 2x - 3$.

Thế vào phương trình thứ nhất: $x + 2(2x - 3) = m$, tức là $5x - 6 = m$, suy ra $x = \dfrac{m + 6}{5}$. Khi đó $y = 2 \cdot \dfrac{m + 6}{5} - 3 = \dfrac{2m - 3}{5}$. Vậy với mọi $m$, hệ có nghiệm duy nhất $(x; y) = \left(\dfrac{m + 6}{5}; \dfrac{2m - 3}{5}\right)$.

Điều kiện $x > 0$ tương đương $\dfrac{m + 6}{5} > 0$; vì $5 > 0$ nên $m + 6 > 0$, tức là $m > -6$.

Điều kiện $y < 0$ tương đương $\dfrac{2m - 3}{5} < 0$, tức là $2m - 3 < 0$, hay $m < \dfrac{3}{2}$.

Hai điều kiện phải thoả mãn đồng thời, nên lấy phần chung của $m > -6$ và $m < \dfrac{3}{2}$: $-6 < m < \dfrac{3}{2}$.

Thử lại với $m = 0$ (thuộc khoảng trên): nghiệm là $\left(\dfrac{6}{5}; -\dfrac{3}{5}\right)$ có $x > 0$ và $y < 0$ (đúng).

Đáp số: $\boxed{-6 < m < \dfrac{3}{2}}$
\end{solution}
\end{problem}

% ---------------------------------------------------------------- đặt ẩn phụ với biểu thức ở mẫu
\begin{problem}[9DS-HPT-046]
% nguon: dummy (tự soạn để mở rộng ngân hàng)
% lop: 9
% kc_du_kien: 9DS-HPT-06, 9DS-HPT-03, 8DS-PS-01
% do_kho_du_kien: 4
Giải hệ phương trình
\[ \begin{cases} \dfrac{2}{x - 1} + \dfrac{4}{y + 2} = 2 \\ \dfrac{6}{x - 1} - \dfrac{8}{y + 2} = 1 \end{cases} \]

\begin{solution}
Điều kiện xác định: các mẫu khác $0$, tức là $x \neq 1$ và $y \neq -2$.

Đặt $u = \dfrac{1}{x - 1}$, $v = \dfrac{1}{y + 2}$, hệ trở thành hệ bậc nhất hai ẩn: $2u + 4v = 2$ và $6u - 8v = 1$.

Nhân hai vế phương trình thứ nhất với $2$ để hệ số của $v$ đối nhau: $4u + 8v = 4$. Cộng vế theo vế với phương trình thứ hai: $10u = 5$, suy ra $u = \dfrac{1}{2}$.

Thế $u = \dfrac{1}{2}$ vào $2u + 4v = 2$: $1 + 4v = 2$, suy ra $v = \dfrac{1}{4}$.

Trở lại ẩn ban đầu: $\dfrac{1}{x - 1} = \dfrac{1}{2}$ nên $x - 1 = 2$, suy ra $x = 3$; $\dfrac{1}{y + 2} = \dfrac{1}{4}$ nên $y + 2 = 4$, suy ra $y = 2$.

Đối chiếu điều kiện: $x = 3 \neq 1$ và $y = 2 \neq -2$ (thoả mãn). Thử lại: $\dfrac{2}{2} + \dfrac{4}{4} = 2$ và $\dfrac{6}{2} - \dfrac{8}{4} = 3 - 2 = 1$ (đúng).

Đáp số: $\boxed{(x; y) = (3; 2)}$
\end{solution}
\end{problem}

% ---------------------------------------------------------------- bài toán hai vòi nước, đặt ẩn phụ
\begin{problem}[9DS-HPT-047]
% nguon: dummy (tự soạn để mở rộng ngân hàng)
% lop: 9
% kc_du_kien: 9DS-HPT-07, 9DS-HPT-06, 9DS-HPT-03
% do_kho_du_kien: 4
Hai vòi nước cùng chảy vào một bể không có nước thì sau $6$ giờ đầy bể. Nếu mở vòi I chảy trong $3$ giờ rồi khoá lại, sau đó mở vòi II chảy trong $6$ giờ thì được $\dfrac{7}{10}$ bể. Hỏi nếu mở riêng từng vòi thì mỗi vòi chảy đầy bể sau bao lâu? Biết lưu lượng của mỗi vòi không đổi.

\begin{solution}
Gọi thời gian vòi I chảy một mình đầy bể là $x$ (giờ), của vòi II là $y$ (giờ); điều kiện $x > 6$, $y > 6$ (một vòi chảy một mình phải lâu hơn khi hai vòi cùng chảy).

Coi cả bể là $1$. Trong một giờ, vòi I chảy được $\dfrac{1}{x}$ bể và vòi II chảy được $\dfrac{1}{y}$ bể.

Hai vòi cùng chảy thì sau $6$ giờ đầy bể, nên trong một giờ cả hai vòi chảy được $\dfrac{1}{6}$ bể: $\dfrac{1}{x} + \dfrac{1}{y} = \dfrac{1}{6}$.

Vòi I chảy $3$ giờ và vòi II chảy $6$ giờ được $\dfrac{7}{10}$ bể: $\dfrac{3}{x} + \dfrac{6}{y} = \dfrac{7}{10}$.

Đặt $u = \dfrac{1}{x}$, $v = \dfrac{1}{y}$ (với $u > 0$, $v > 0$), ta được hệ bậc nhất hai ẩn: $u + v = \dfrac{1}{6}$ và $3u + 6v = \dfrac{7}{10}$.

Nhân hai vế phương trình thứ nhất với $3$ để hệ số của $u$ bằng nhau: $3u + 3v = \dfrac{1}{2}$. Trừ vế theo vế phương trình $3u + 6v = \dfrac{7}{10}$ cho phương trình này: $3v = \dfrac{7}{10} - \dfrac{1}{2} = \dfrac{1}{5}$, suy ra $v = \dfrac{1}{15}$.

Thế $v = \dfrac{1}{15}$ vào phương trình thứ nhất: $u = \dfrac{1}{6} - \dfrac{1}{15} = \dfrac{5 - 2}{30} = \dfrac{1}{10}$.

Trở lại ẩn ban đầu: $\dfrac{1}{x} = \dfrac{1}{10}$ nên $x = 10$; $\dfrac{1}{y} = \dfrac{1}{15}$ nên $y = 15$.

Đối chiếu điều kiện: $x = 10 > 6$, $y = 15 > 6$ (thoả mãn). Thử lại: $\dfrac{1}{10} + \dfrac{1}{15} = \dfrac{3 + 2}{30} = \dfrac{1}{6}$ và $\dfrac{3}{10} + \dfrac{6}{15} = \dfrac{3}{10} + \dfrac{4}{10} = \dfrac{7}{10}$ (đúng).

Đáp số: $\boxed{\text{Vòi I: } 10 \text{ giờ; vòi II: } 15 \text{ giờ}}$
\end{solution}
\end{problem}

% ---------------------------------------------------------------- hệ thức không phụ thuộc tham số, giá trị nhỏ nhất của tích
\begin{problem}[9DS-HPT-048]
% nguon: dummy (tự soạn để mở rộng ngân hàng)
% lop: 9
% kc_du_kien: 9DS-HPT-09, 9DS-HPT-03, 8DS-HDT-06
% do_kho_du_kien: 4
Cho hệ phương trình (với $m$ là tham số)
\[ \begin{cases} x + y = 2m - 2 \\ 2x - y = m + 5 \end{cases} \]
a) Giải hệ phương trình theo $m$.

b) Tìm một hệ thức liên hệ giữa $x$ và $y$ không phụ thuộc vào $m$.

c) Tìm giá trị nhỏ nhất của tích $P = xy$ và giá trị của $m$ tương ứng.

\begin{solution}
a) Cộng vế theo vế hai phương trình để khử $y$: $3x = 3m + 3$, suy ra $x = m + 1$.

a) Thế $x = m + 1$ vào phương trình thứ nhất: $y = 2m - 2 - (m + 1) = m - 3$. Vậy với mọi $m$, hệ có nghiệm duy nhất $(x; y) = (m + 1; m - 3)$.

a) Thử lại phương trình thứ hai: $2(m + 1) - (m - 3) = m + 5$ (đúng).

b) Trừ vế theo vế hai biểu thức của nghiệm để khử $m$: $x - y = (m + 1) - (m - 3) = 4$. Vậy $x - y = 4$ với mọi $m$.

c) Thay nghiệm vào tích: $P = (m + 1)(m - 3) = m^2 - 2m - 3$.

c) Theo hằng đẳng thức bình phương của một hiệu: $P = (m^2 - 2m + 1) - 4 = (m - 1)^2 - 4$.

c) Vì $(m - 1)^2 \geq 0$ với mọi $m$ nên $P \geq -4$. Dấu bằng xảy ra khi $(m - 1)^2 = 0$, tức là $m = 1$.

c) Thử lại: với $m = 1$ nghiệm là $(2; -2)$ và $P = 2 \cdot (-2) = -4$ (đúng).

Đáp số: a) $\boxed{(x; y) = (m + 1; m - 3)}$; b) $\boxed{x - y = 4}$; c) $\boxed{\min P = -4 \text{ khi } m = 1}$
\end{solution}
\end{problem}

% ---------------------------------------------------------------- giao điểm hai đường thẳng, diện tích tam giác
\begin{problem}[9DS-HPT-049]
% nguon: dummy (tự soạn để mở rộng ngân hàng)
% lop: 9
% kc_du_kien: 9DS-HPT-02, 8DS-HSBN-06, 8DS-HSBN-03
% do_kho_du_kien: 4
Cho hai đường thẳng $(d_1)\colon y = x + 2$ và $(d_2)\colon y = -2x + 8$.

a) Tìm toạ độ giao điểm $A$ của $(d_1)$ và $(d_2)$.

b) Gọi $B$ và $C$ lần lượt là giao điểm của $(d_1)$ và $(d_2)$ với trục hoành. Tính diện tích tam giác $ABC$ (đơn vị đo trên các trục toạ độ là xentimét).

\begin{solution}
a) Toạ độ giao điểm $A$ là nghiệm của hệ $y = x + 2$ và $y = -2x + 8$. Thế $y = x + 2$ vào phương trình thứ hai: $x + 2 = -2x + 8$.

a) Chuyển vế: $3x = 6$, suy ra $x = 2$; khi đó $y = 2 + 2 = 4$. Vậy $A(2; 4)$.

a) Thử lại với $(d_2)$: $-2 \cdot 2 + 8 = 4$ (đúng).

b) Điểm $B$ thuộc trục hoành nên có tung độ $0$. Thay $y = 0$ vào $(d_1)$: $0 = x + 2$, suy ra $x = -2$, nên $B(-2; 0)$.

b) Tương tự, thay $y = 0$ vào $(d_2)$: $0 = -2x + 8$, suy ra $x = 4$, nên $C(4; 0)$.

b) Vì $B$, $C$ cùng nằm trên trục hoành nên $BC = |4 - (-2)| = 6$ (cm). Đường cao của tam giác $ABC$ hạ từ $A$ xuống trục hoành bằng tung độ của $A$, tức là $4$ cm.

b) Diện tích tam giác $ABC$: $S = \dfrac{1}{2} \cdot BC \cdot 4 = \dfrac{1}{2} \cdot 6 \cdot 4 = 12$ (cm$^2$).

Đáp số: a) $\boxed{A(2; 4)}$; b) $\boxed{S = 12\ \text{cm}^2}$
\end{solution}
\end{problem}

% ---------------------------------------------------------------- biện luận ba trường hợp, giá trị nguyên
\begin{problem}[9DS-HPT-050]
% nguon: dummy (tự soạn để mở rộng ngân hàng)
% lop: 9
% kc_du_kien: 9DS-HPT-05, 9DS-HPT-04, 9DS-HPT-02, 8DS-PTNT-02
% do_kho_du_kien: 5
Cho hệ phương trình (với $m$ là tham số)
\[ \begin{cases} x + my = m + 1 \\ mx + 4y = 2m + 2 \end{cases} \]
a) Tìm $m$ để hệ có nghiệm duy nhất và tìm nghiệm đó theo $m$.

b) Khi $m = 2$ và khi $m = -2$, hệ có bao nhiêu nghiệm?

c) Tìm các giá trị nguyên của $m$ để hệ có nghiệm duy nhất $(x; y)$ mà $x + y$ là số nguyên.

\begin{solution}
a) Từ phương trình thứ nhất rút $x$ theo $y$: $x = m + 1 - my$.

a) Thế vào phương trình thứ hai: $m(m + 1 - my) + 4y = 2m + 2$, tức là $m^2 + m - m^2y + 4y = 2m + 2$; chuyển vế và nhóm các hạng tử chứa $y$ được $(4 - m^2)y = -m^2 + m + 2$.

a) Phân tích hai vế thành nhân tử (vế trái dùng hằng đẳng thức hiệu hai bình phương, vế phải dùng $-m^2 + m + 2 = (2 - m)(m + 1)$): $(2 - m)(2 + m)y = (2 - m)(m + 1)$. Hệ đã cho tương đương với hệ gồm phương trình này và $x = m + 1 - my$.

a) Hệ có nghiệm duy nhất khi và chỉ khi phương trình $(2 - m)(2 + m)y = (2 - m)(m + 1)$ có nghiệm duy nhất, tức là $(2 - m)(2 + m) \neq 0$, hay $m \neq 2$ và $m \neq -2$.

a) Với $m \neq \pm 2$: chia hai vế cho $(2 - m)(2 + m)$ được $y = \dfrac{m + 1}{m + 2}$; khi đó $x = m + 1 - \dfrac{m(m + 1)}{m + 2} = \dfrac{(m + 1)(m + 2 - m)}{m + 2} = \dfrac{2(m + 1)}{m + 2}$.

b) Khi $m = 2$, phương trình $(2 - m)(2 + m)y = (2 - m)(m + 1)$ thành $0y = 0$, đúng với mọi $y$. Hệ có vô số nghiệm: $x = 3 - 2y$ với $y$ tuỳ ý.

b) Khi $m = -2$, phương trình đó thành $0y = 4 \cdot (-1) = -4$, không có $y$ nào thoả mãn. Hệ vô nghiệm.

c) Với $m \neq \pm 2$: $x + y = \dfrac{2(m + 1) + (m + 1)}{m + 2} = \dfrac{3(m + 1)}{m + 2} = \dfrac{3(m + 2) - 3}{m + 2} = 3 - \dfrac{3}{m + 2}$.

c) Với $m$ nguyên, $x + y$ là số nguyên khi và chỉ khi $\dfrac{3}{m + 2}$ là số nguyên, tức là số nguyên $m + 2$ là ước của $3$: $m + 2 \in \{-3;\ -1;\ 1;\ 3\}$.

c) Suy ra $m \in \{-5;\ -3;\ -1;\ 1\}$; cả bốn giá trị đều khác $2$ và $-2$ (thoả mãn điều kiện có nghiệm duy nhất).

c) Thử lại: $m = -5$ cho $x + y = 3 - \dfrac{3}{-3} = 4$; $m = -3$ cho $3 - \dfrac{3}{-1} = 6$; $m = -1$ cho $3 - 3 = 0$; $m = 1$ cho $3 - 1 = 2$ (đều là số nguyên).

Đáp số: a) $\boxed{m \neq \pm 2,\ (x; y) = \left(\dfrac{2(m + 1)}{m + 2}; \dfrac{m + 1}{m + 2}\right)}$; b) $\boxed{m = 2 \text{: vô số nghiệm; } m = -2 \text{: vô nghiệm}}$; c) $\boxed{m \in \{-5;\ -3;\ -1;\ 1\}}$
\end{solution}
\end{problem}

% ---------------------------------------------------------------- đặt ẩn phụ hai lần
\begin{problem}[9DS-HPT-051]
% nguon: dummy (tự soạn để mở rộng ngân hàng)
% lop: 9
% kc_du_kien: 9DS-HPT-06, 9DS-HPT-03, 8DS-PS-01
% do_kho_du_kien: 5
Giải hệ phương trình
\[ \begin{cases} \dfrac{3}{2x - y} + \dfrac{2}{x + 2y} = -1 \\ \dfrac{6}{2x - y} - \dfrac{1}{x + 2y} = 3 \end{cases} \]

\begin{solution}
Điều kiện xác định: các mẫu khác $0$, tức là $2x - y \neq 0$ và $x + 2y \neq 0$.

Đặt $u = \dfrac{1}{2x - y}$, $v = \dfrac{1}{x + 2y}$, hệ trở thành hệ bậc nhất hai ẩn: $3u + 2v = -1$ và $6u - v = 3$.

Từ phương trình thứ hai rút $v$ theo $u$: $v = 6u - 3$.

Thế vào phương trình thứ nhất: $3u + 2(6u - 3) = -1$, tức là $15u - 6 = -1$, suy ra $u = \dfrac{1}{3}$; khi đó $v = 6 \cdot \dfrac{1}{3} - 3 = -1$.

Trở lại biểu thức ban đầu: $\dfrac{1}{2x - y} = \dfrac{1}{3}$ nên $2x - y = 3$; $\dfrac{1}{x + 2y} = -1$ nên $x + 2y = -1$. Đây là một hệ bậc nhất hai ẩn $x$, $y$ mới.

Nhân hai vế phương trình $2x - y = 3$ với $2$ để hệ số của $y$ đối nhau: $4x - 2y = 6$. Cộng vế theo vế với $x + 2y = -1$: $5x = 5$, suy ra $x = 1$; khi đó $y = 2x - 3 = -1$.

Đối chiếu điều kiện: $2 \cdot 1 - (-1) = 3 \neq 0$ và $1 + 2 \cdot (-1) = -1 \neq 0$ (thoả mãn). Thử lại: $\dfrac{3}{3} + \dfrac{2}{-1} = -1$ và $\dfrac{6}{3} - \dfrac{1}{-1} = 3$ (đúng).

Đáp số: $\boxed{(x; y) = (1; -1)}$
\end{solution}
\end{problem}

\end{document}
